%2multibyte Version: 5.50.0.2960 CodePage: 1252

\documentclass[12pt]{article}
%%%%%%%%%%%%%%%%%%%%%%%%%%%%%%%%%%%%%%%%%%%%%%%%%%%%%%%%%%%%%%%%%%%%%%%%%%%%%%%%%%%%%%%%%%%%%%%%%%%%%%%%%%%%%%%%%%%%%%%%%%%%%%%%%%%%%%%%%%%%%%%%%%%%%%%%%%%%%%%%%%%%%%%%%%%%%%%%%%%%%%%%%%%%%%%%%%%%%%%%%%%%%%%%%%%%%%%%%%%%%%%%%%%%%%%%%%%%%%%%%%%%%%%%%%%%
\usepackage{amsfonts}
\usepackage{amsmath}

\setcounter{MaxMatrixCols}{10}
%TCIDATA{OutputFilter=LATEX.DLL}
%TCIDATA{Version=5.50.0.2960}
%TCIDATA{Codepage=1252}
%TCIDATA{<META NAME="SaveForMode" CONTENT="1">}
%TCIDATA{BibliographyScheme=Manual}
%TCIDATA{Created=Monday, September 22, 2003 09:40:11}
%TCIDATA{LastRevised=Monday, September 29, 2025 17:17:27}
%TCIDATA{<META NAME="GraphicsSave" CONTENT="32">}
%TCIDATA{<META NAME="DocumentShell" CONTENT="General\SW\Blank - Standard LaTeX Article">}
%TCIDATA{Language=American English}
%TCIDATA{CSTFile=LaTeX article (bright).cst}

\newtheorem{theorem}{Theorem}
\newtheorem{acknowledgement}[theorem]{Acknowledgement}
\newtheorem{algorithm}[theorem]{Algorithm}
\newtheorem{axiom}[theorem]{Axiom}
\newtheorem{case}[theorem]{Case}
\newtheorem{claim}[theorem]{Claim}
\newtheorem{conclusion}[theorem]{Conclusion}
\newtheorem{condition}[theorem]{Condition}
\newtheorem{conjecture}[theorem]{Conjecture}
\newtheorem{corollary}[theorem]{Corollary}
\newtheorem{criterion}[theorem]{Criterion}
\newtheorem{definition}[theorem]{Definition}
\newtheorem{example}[theorem]{Example}
\newtheorem{exercise}[theorem]{Exercise}
\newtheorem{lemma}[theorem]{Lemma}
\newtheorem{notation}[theorem]{Notation}
\newtheorem{problem}[theorem]{Problem}
\newtheorem{proposition}[theorem]{Proposition}
\newtheorem{remark}[theorem]{Remark}
\newtheorem{solution}[theorem]{Solution}
\newtheorem{summary}[theorem]{Summary}
\newenvironment{proof}[1][Proof]{\textbf{#1.} }{\ \rule{0.5em}{0.5em}}
\input{tcilatex}
\input{tcilatex}
\setlength{\textwidth}{6.5in}
\setlength{\textheight}{9in}
\setlength{\topmargin}{-.50in}
\setlength{\oddsidemargin}{.0in}
\setlength{\evensidemargin}{.0in}

\begin{document}


\begin{center}
{\LARGE Assignment 5}
\end{center}

\noindent Blankenau\hspace*{\stretch{0.7500}}

\noindent Economics 905, Fall 2025

\noindent \textit{Instructions.} You may work in groups and may compare
answers prior to submission. Joint submissions are allowed and encouraged.
All work should be legible and concise. A due date will be announced in
class.\medskip 

\begin{enumerate}
\item Consider an economy consisting of overlapping generations of
two-period-lived homogeneous agents, a representative firm producing a
single good and a government. A representative agent is born each period and
the population of each generation is normalized to 1. The initial young and
old agents are endowed with human capital $h_{0}.$ All agents have one unit
of time to supply exogenously to the labor market in each period of life.
Agents beyond the initial period receive an endowment of human capital which
depends on government and family spending on education as the first period
of life begins. Specifically, the human capital of the representative agent
born in period $t$ is $h_{t}=\left( f_{t}+g_{t}\right) ^{\alpha }$ where $%
f_{t}$ is spending by the agent's parents on education in period $t$ and $%
g_{t}$ is spending by the government\ on education in period $t$. Assume $%
0\leq \alpha \leq 1.$ In the second period of life, through experience, the
agent's level of human capital grows to $\gamma h_{t}$ where $\gamma >1.$
Parents value consumption in each period of life. They also value the human
capital of their children. Preferences are given by 
\begin{equation*}
U_{t}=\ln c_{t,t}+\beta \ln c_{t,t+1}+\psi \ln h_{t+1}
\end{equation*}%
where $U_{t}$ is the lifetime utility of an agent born in period $t$, $%
c_{t,t}$ and $c_{t,t+1}$ are this agent's consumption when young and old and 
$h_{t+1}$ is the human capital of the agent's children. There is no storage
or capital in this economy so the structure precludes borrowing and lending.
Output is produced according to $Y_{t}=AH_{t}$ where $H_{t}$ is the amount
of human capital employed by the firm and $A>0$ is a scalar. Each period $t$
agent inelastically provides $h_{t}$ units of human capital per unit of time
when young and $\gamma h_{t}$ when old so that income is $\omega _{t}h_{t}$
and $\omega _{t+1}\gamma h_{t}$ when young and old where $\omega _{t}$ is
the wage per unit of human capital in period $t$. Government lump sum taxes
the old to pay for its education expenditure subject to a balanced budget.

\begin{enumerate}
\item Write down the problem for a period $t$ agent and find the first order
condition. The agent takes $h_{t}$, $\omega _{t},$ $\omega _{t+1}$ and the
tax rate $\tau _{t,t+1}$ as given. Also write down the firm's problem and
give the government's budget constraint.

\item Solve for the steady state level of $f$ in this economy. How does
government education spending affect total education spending?

\item Now suppose the situation is the same except for two differences.
First, government lump sum taxes the young to pay for its education
expenditure so that $c_{t,t}=\left( \omega _{t}h_{t}-\tau _{t,t}\right) $
and $c_{t,t+1}=\left( \omega _{t+1}h_{t}\gamma -f_{t+1}\right) $. Second,
the human capital accumulation function is $h_{t}=f_{t}^{\alpha
}g_{t}^{1-\alpha }$. Write down the problem for a period $t$ agent and find
the first order condition. The agent takes $h_{t}$, $\omega _{t},$ $\omega
_{t+1}$ and the tax rate $\tau _{t,t}$ as given. Also write down the firm's
problem and give the government's budget constraint. Find the steady state
level of $f$ in this economy.
\end{enumerate}
\end{enumerate}

\pagebreak

\begin{enumerate}
\item 
\begin{enumerate}
\item ...%
\begin{equation*}
\underset{c_{t,t},c_{t,t+1},f_{t+1}}{\max }\ln c_{t,t}+\beta \ln
c_{t,t+1}+\psi \ln h_{t+1}
\end{equation*}%
subject to%
\begin{eqnarray*}
c_{t,t} &=&\omega _{t}h_{t} \\
c_{t,t+1} &=&\omega _{t+1}h_{t}\gamma -\tau _{t,t+1}-f_{t+1} \\
h_{t+1} &=&E\left( f_{t+1}+g_{t+1}\right) ^{\alpha }
\end{eqnarray*}%
or 
\begin{equation*}
\underset{,f_{t+1}}{\max }\ln \left( \omega _{t}h_{t}\right) +\beta \ln
\left( \omega _{t+1}h_{t}\gamma -\tau _{t,t+1}-f_{t+1}\right) +\psi \ln
\left( f_{t+1}+g_{t+1}\right) ^{\alpha }
\end{equation*}%
so the foc is 
\begin{equation*}
\frac{\beta }{w_{t+1}h_{t}\gamma -\tau _{t,t+1}-f_{t+1}}=\frac{\alpha \psi }{%
f_{t+1}+g_{t+1}}
\end{equation*}%
Firm:%
\begin{equation*}
\underset{H_{t}}{\max }AH_{t}-\omega _{t}H_{t}
\end{equation*}%
Government%
\begin{equation*}
g_{t+1}=\tau _{t,t+1}
\end{equation*}

\item By prior arguments, the $\omega _{t}=A$ for all time periods. Our
first order condition becomes 
\begin{equation*}
\frac{\beta }{Ah\gamma -\tau -f}=\frac{\alpha \psi }{f+g}
\end{equation*}%
or 
\begin{equation*}
\frac{\beta }{A\left( f+g\right) ^{\alpha }\gamma -g-f}=\frac{\alpha \psi }{%
f+g}
\end{equation*}%
\begin{equation*}
\beta \left( f+g\right) =\alpha \psi \left( A\left( f+g\right) ^{\alpha
}\gamma -\left( f+g\right) \right)
\end{equation*}%
\begin{equation*}
\beta =\alpha \psi \left( A\left( f+g\right) ^{\alpha -1}\gamma -1\right)
\end{equation*}%
\begin{equation*}
\frac{\beta +\alpha \psi }{\alpha \psi A\gamma }=\left( f+g\right) ^{\alpha
-1}
\end{equation*}%
\begin{equation*}
f=\left( \frac{\alpha \psi A\gamma }{\beta +\alpha \psi }\right) ^{\frac{1}{%
1-\alpha }}-g
\end{equation*}%
Since $g$ cannot be less than 0%
\begin{equation*}
f=\max \left\{ \left( \frac{\alpha \psi A\gamma }{\beta +\alpha \psi }%
\right) ^{\frac{1}{1-\alpha }}-g,0\right\}
\end{equation*}

\item 
\begin{equation*}
\underset{f_{t+1}}{\max }\ln \left( \omega _{t}h_{t}-\tau _{t,t}\right)
+\beta \ln \left( \omega _{t+1}h_{t}\gamma -f_{t+1}\right) +\psi \ln \left(
f_{t+1}^{\alpha }g_{t+1}^{1-\alpha }\right) 
\end{equation*}%
\begin{equation*}
g_{t+1}=\tau _{t,t}
\end{equation*}%
firm%
\begin{equation*}
\underset{H_{t}}{\max }AH_{t}-\omega _{t}H_{t}
\end{equation*}%
\begin{equation*}
\frac{\beta }{\left( \omega _{t+1}h_{t}\gamma -f_{t+1}\right) }+\alpha \psi 
\frac{f_{t+1}^{\alpha -1}g_{t+1}^{1-\alpha }}{f_{t+1}^{\alpha
}g_{t+1}^{1-\alpha }}
\end{equation*}%
In a steady state%
\begin{equation*}
\frac{\beta }{Af^{\alpha }g^{1-\alpha }\gamma -f}=\frac{\alpha \psi }{f}
\end{equation*}%
from which 
\begin{equation*}
\beta f=\alpha \psi \left( Af^{\alpha }g^{1-\alpha }\gamma -f\right) 
\end{equation*}%
and solving gives%
\begin{equation*}
\left( \frac{A\alpha \psi \gamma }{\beta +\alpha \psi }\right) ^{\frac{1}{%
1-\alpha }}g=f.
\end{equation*}
\end{enumerate}
\end{enumerate}

\end{document}
